\chapter*{Introduction}
\addcontentsline{toc}{chapter}{Introduction}
This thesis stems from the question at the core of the project I presented more than four years ago to apply to the PhD program: "Would an Universal Basic Income have inflationary effects?".

In these four years the answer I came to is "It doesn't matter", but the reasoning behind it doesn't belong to what is part of Economics today (a hint is given in the chapter \ref{ch:against}).
Rather, this thesis explore how an economist can answer a similar question from a methodological point of view in the first part and going deeper into discussing how agent-based models can help (and their limits and upsides) in the second part.
The model presented is not sufficient to answer the original question, having too few details, but is presented as a teaching tool to explore how Stock-Flow Consistent models can be expanded into agent-based ones, focusing on production.
Developing a model from scratch requires, as many older scholars told me, to adopt the KISS approach: Keep It Simple Stupid. And the best way to keep something simple is adding a block at time, but, as I'll show, each block add complexity and makes the model unstable, very reactive to small changes in parameter (and so very difficult to calibrate on realistic data, if at all possible).

The model presented in part \ref{part:model} shows some interesting, and probably innovative, traits, focusing on the production side of the economy rather than financial one (as commonly done in SFC models). But I believe it's impact is more significative as a introductory and teaching tool rather than a research one.

Chapter \ref{ch:complexity} is written as a scientific paper, and currently looking for its own publication, while the subsequent chapter address its shortcomings and bring further reflections on the topic of complexity anf knowability and modellability of reality from a post-modernist perspective.
